% 分野別類題: ガウス積分
% 2重積分＋極座標変換、パラメータ微分、平方完成の3問。
% コンパイル: TEXINPUTS="../../../00_common:$TEXINPUTS" latexmk -xelatex problems.tex
\documentclass[a4paper,11pt]{article}
\usepackage{preamble}

\begin{document}

\kamokumei{類題演習 --- ガウス積分}

\shijibun{以下の各問に答えよ。導出過程を明示すること。（ヒント: 2重積分を極座標変換すること、あるいは積分をパラメータの関数とみて微分することで計算できる場合がある。）}

\shomon{問1.}{次のガウス積分の値を、2重積分と極座標変換を用いて導出せよ。}
\[
I = \int_{-\infty}^{\infty} e^{-x^{2}}\, dx
\]
さらに、これを用いて $a > 0$ に対する
\[
\int_{-\infty}^{\infty} e^{-a x^{2}}\, dx
\]
の値を求めよ。

\shomon{問2.}{$a > 0$ とする。次の積分の値を求めよ。}
\[
\text{(1)}\quad \int_{0}^{\infty} x\, e^{-a x^{2}}\, dx
\qquad\qquad
\text{(2)}\quad \int_{-\infty}^{\infty} x^{2}\, e^{-a x^{2}}\, dx
\]
（ヒント: (2) は $\displaystyle\int_{-\infty}^{\infty} e^{-a x^{2}}\, dx$ を $a$ の関数とみて、$a$ で微分することで求められる。）

\shomon{問3.}{$a > 0$、$b$ を実数の定数とする。平方完成を用いて次の積分の値を求めよ。}
\[
\int_{-\infty}^{\infty} e^{-a x^{2} + b x}\, dx
\]
また、この結果を用いて
\[
\int_{-\infty}^{\infty} e^{-x^{2} + 2x}\, dx
\]
の値を求めよ。

\end{document}
